
\documentclass[12pt,aspectratio=169,t,fleqn]{beamer}

\usepackage[fontset=none,autoindent=true]{ctex}
\IfFontExistsTF{Songti SC}{%
	\setCJKmainfont{Songti SC}[AutoFakeSlant=0.2]%
	\setCJKsansfont{Songti SC}[AutoFakeSlant=0.2]%
	\setCJKmonofont{Songti SC}%
}{%
	\setCJKmainfont{FandolSong-Regular}[Extension=.otf,BoldFont=FandolSong-Bold,AutoFakeSlant=0.2]%
	\setCJKsansfont{FandolSong-Regular}[Extension=.otf,BoldFont=FandolSong-Bold,AutoFakeSlant=0.2]%
	\setCJKmonofont{FandolFang-Regular}[Extension=.otf,BoldFont=FandolFang-Regular]%
}
\usepackage{hyperref}
\setmainfont{XCharter} % Nicer fonts — the charter package is ignored under XeLaTeX, so use fontspec instead
% 等宽字体：优先 Maple Mono NF（静态字族，真实粗/斜体），其次可变版，未安装则回退默认
\IfFontExistsTF{Maple Mono NF}{%
	\setmonofont{Maple Mono NF}[Scale=MatchLowercase,Ligatures=TeX]%
}{%
	\IfFontExistsTF{Maple Mono}{%
		\setmonofont{Maple Mono}[Scale=MatchLowercase,Ligatures=TeX]%
	}{}%
}
\usefonttheme{professionalfonts}%**强烈建议加上**，关闭 beamer 自带字体替换，防止字体冲突，数学公式字体正常渲染
\mode<presentation>
{
  \usetheme{Warsaw}

  \setbeamercovered{invisible}
 %   含义：**大括号内的代码，只在 presentation（幻灯片）模式下生效；
 %   如果切换成 article 讲义模式，这一段直接忽略不执行。**
 
% 控制**被覆盖（还没出现）的条目**样式： 
% `transparent`：未出现的条目变浅、半透明（灰色淡化，占位还在页面上）
% 可选：`invisible`（默认，没出现的内容直接隐藏，页面空白）
% 对应你之前的 `<+->` 逐条动画：
% `transparent`：后面还没出来的 \item 是淡淡的灰色，提前看到占位
% `invisible`：没到显示时机的 item 完全看不见
}


\title[2.1等式性质与不等式性质(习题课)] % (optional, use only with long paper titles)
{2.1等式性质与不等式性质(习题课)}

\subtitle
{}

\author[好好学习，努力向上] % (optional, use only with lots of authors)
{
%	金永权\inst{1} \and 某某某\inst{2}
	}

\institute[夏邑一高] % (optional, but mostly needed)
{
%  \inst{1}%
%  高一年级
%  22 班
%  \and
%  \inst{2}%
% 高一年级
% 24班
 }
% - Use the \inst command only if there are several affiliations.
% - Keep it simple, no one is interested in your street address.

\date[2026.9.19] % (optional, should be abbreviation of conference name)
{2026.9.19}
% - Either use conference name or its abbreviation.
% - Not really informative to the audience, more for people (including
%   yourself) who are reading the slides online

\subject{数学科}
% This is only inserted into the PDF information catalog. Can be left out. 


\AtBeginSubsection[]
{
  \begin{frame}<beamer>{主要内容}
    \tableofcontents[sectionstyle=show/hide,subsectionstyle=show/shaded]
  \end{frame}
}
%	作用：**每次开启一个 `\subsection` 小节时，自动执行大括号里面的代码**。
%	`[]` 里面可以放 overlay 规范，空括号代表默认。
%	含义：**这个目录帧只在 presentation（幻灯片模式）编译出来；
%	如果切换成 article 讲义模式，这一帧直接忽略，不会生成**。
%	等价写法：`\mode<presentation>{\begin{frame}...\end{frame}}`
%	目录命令的选项：

%	- `currentsection`：**只高亮当前所在的大 section，其他 section 变灰**
%	- `currentsubsection`：**高亮当前 subsection，同 section 下其他小节变浅**
%\beamerdefaultoverlayspecification{<+->}
\setbeamerfont{itemize/enumerate body}{shape=\upshape}    % 正文直立
\setbeamerfont{proof body}{shape=\upshape}
\setbeamerfont{solution body}{shape=\upshape}
\begin{document}

\begin{frame}
  \titlepage
\end{frame}

\begin{frame}{主要内容}
  \tableofcontents
  % You might wish to add the option [pausesections]
\end{frame}


\section{课本习题}

\subsection{P43第5题}

\begin{frame}{P43第5题}
  % - A title should summarize the slide in an understandable fashion
  %   for anyone how does not follow everything on the slide itself.

  \begin{solution}
  	$\because\;2<a<3,-2<b<-1$， 
  	
  	$\therefore\;4<2a<6,2<2a+b<5$。
  \end{solution}
  
\end{frame}


\subsection{P43第7题}

\begin{frame}{P43第7题}

  \begin{proof}
  	$\because\;c<d<0$，
  	
  	$\therefore\;-c>-d>0$，
  	
  	$\because\;a>b>0$，
  	
  	$\therefore\;a-c>b-d>0$，$\frac{1}{a-c}<\frac{1}{b-d}$，
  	
  	$\because\;e<0$，
  	
  	$\therefore\;\frac{e}{a-c}>\frac{e}{b-d}$。
  \end{proof}
\end{frame}

\subsection{P43第8题}
\begin{frame}{P43第8题}
\begin{solution}
	\begin{enumerate}[(\Alph*)]
		\item 若$c=0$，结论不成立；
		\item 根据不等式性质6，结论成立；
		\item 因为$a<b<0$，所以$a^2 > ab>0$，$ab>b^2>0$，因此$a^2>ab>b^2$。
		\item 因为$a<b<0$，所以$ab>0,\frac{1}{ab}>0$，$a\cdot\frac{1}{ab}<b\cdot\frac{1}{ab}$，
		即\[
		\frac{1}{b} < \frac{1}{a},\text{或} \frac{1}{a} > \frac{1}{b}
		\]
	\end{enumerate}
\end{solution}
\end{frame}

\subsection{P43第9题}
\begin{frame}{P43第9题}
	\begin{solution}
		设圆的半径为$r$，正方形的边长为$a$，由$2\pi r = 4a$，得$r=\frac{2a}{\pi}$，所以$\pi r^2 -a^2 = \frac{(4-\pi)a^2}{\pi}>0$。所以周长相等的圆和正方形，圆的面积大于正方形的面积。
	\end{solution}
\end{frame}

\subsection{P43第10题}
\begin{frame}{P43第10题}
	\begin{solution}
	\upshape	所列不等式为\[
		\frac{a}{b}<\frac{m+a}{m+b}
		\]
		因为$b>a$，所以$\frac{a}{b} - \frac{m+a}{m+b} = \frac{(a-b)m}{(m+b)b} < 0$，
		
		于是，$\frac{a}{b} < \frac{m+a}{m+b}$。
	\end{solution}
\end{frame}

\section{课本习题--拓广探索}

\subsection{P43第11题}
\begin{frame}{P43第11题}
	\begin{proof}
		因为$a>b>0$，所以$\sqrt{a} - \sqrt{b} = \frac{a-b}{\sqrt{a} + \sqrt{b}} >0$，
		
		于是，$\sqrt{a} > \sqrt{b}$。
	\end{proof}
\end{frame}

\subsection{P43第12题}
\begin{frame}{P43第12题}
	\begin{solution}
		\upshape
		设安排甲种货箱 $x$ 节，乙种货箱 $y$ 节，总运费为 $z$。根据题意，得
	
		\begin{align}
			35x+25y &\ge 1530, \tag{1}\\
			15x+35y &\ge 1150, \tag{2}\\
			x+y &= 50. \tag{3}
		\end{align}
		
		
		
	\end{solution}
\end{frame}


\begin{frame}{P43第12题}
	\begin{solution}[接上页]
		\upshape
		把(3)代入(1)，并解得 $x\ge 28$；把(3)代入(2)，并解得 $x\le 30$。于是
		\[
		\begin{cases}
			x=28,\\
			y=22.
		\end{cases}
		\quad\text{或}\quad
		\begin{cases}
			x=29,\\
			y=21.
		\end{cases}
		\quad\text{或}\quad
		\begin{cases}
			x=30,\\
			y=20.
		\end{cases}
		\]
		
		因此，共有三种方案：方案一，安排甲种货箱28节，乙种货箱22节；
		
		方案二，安排甲种货箱29节，乙种货箱21节；
		
		方案三，安排甲种货箱30节，乙种货箱20节。
		
		当甲货箱30节，乙货箱20节时，总运费最少，运费为 $0.5\times 30+0.8\times 20=31$（万元）。
	\end{solution}
\end{frame}

\begin{frame}{}
	
\end{frame}
\end{document}


